\documentclass[10pt,a4paper]{amsart}
\usepackage[margin=19mm]{geometry}
\usepackage{amsmath,amssymb,mathtools}
\usepackage[scr=rsfs,cal=euler]{mathalfa}
\usepackage{xfrac}
\usepackage[unicode,hidelinks,bookmarks=false]{hyperref}
\newcommand{\ie}{i.e.\ }
\newcommand{\cf}{cf.\ }
\newcommand{\resp}{respectively\ }
\newcommand{\Subsection}{\subsection}
\providecommand{\nobreakdash}{\nobreak\hbox{-}}
\setlength{\emergencystretch}{2em}
\multlinegap0pt
\usepackage{amssymb}
\usepackage{xcolor}
\usepackage{enumerate}
\usepackage{dsfont}
\newtheorem{lemma}{Lemma}
\newcommand{\BC}{\mathbb{C}}
\newcommand{\BF}{\mathbb{F}}
\newcommand{\BK}{\mathbb{K}}
\newcommand{\GL}{\mathrm{GL}}
\newcommand{\cb}[1]{\left\{#1\right\}}
\newcommand{\fk}[1]{\left[#1\right]}
\newcommand{\jk}[1]{\left\langle#1\right\rangle}
\newcommand{\jz}[4]{{\begin{pmatrix}
		#1&#2\\
		#3&#4
\end{pmatrix}}}
\newcommand{\lcm}{\mathrm{lcm}}
\newcommand{\ls}{\leqslant}
\newcommand{\one}{\mathds{1}}
\newcommand{\pt}[1]{\left(#1\right)}
\title{A\lowercase{n} E\lowercase{rd\H{o}s}-K\lowercase{o}-R\lowercase{ado result for some principal series representations}}
\author{Jiaqi Liao, Guiying Yan}
\date{Source: arXiv:2604.01953v2 (CC BY 4.0)}
\begin{document}
\maketitle

\noindent\emph{Source and licence.} A\lowercase{n} E\lowercase{rd\H{o}s}-K\lowercase{o}-R\lowercase{ado result for some principal series representations}. Version arXiv:2604.01953v2 (CC BY 4.0), distributed by arXiv under the Creative Commons Attribution 4.0 International licence, http://creativecommons.org/licenses/by/4.0/, as recorded in the arXiv licence metadata for this exact version and shown on the abstract page https://arxiv.org/abs/2604.01953v2. Full licence notice: \texttt{licenses/arxiv-2604.01953-CC-BY-4.0.txt}.\par\smallskip

\noindent\textit{Editorial scope.} Counted unit: the article's lemma eta\_lemma (source0.tex lines 370--380). The article's complete original proof of it is delivered with it (source0.tex lines 382--448, one \texttt{proof} environment of 67 lines). No mathematics is written, repaired or re-derived: the statement and the proof are the article's own lines, in the article's own notation, and no retained line is substituted or re-flowed.  No further source-local context is needed: every cross-reference of every retained line resolves inside the two delivered spans.  Nothing was omitted from any retained span, so the package carries no omission record.  No retained line contains a citation, so the wrapper carries no bibliography and no bibliography entry.  No retained line contains a figure environment, an image-inclusion command or drawing code, so no image is delivered and the package declares no asset dependency.  Editorial formatting only, in addition: an AMS \texttt{amsart} preamble that restates the article's own declarations and macros used by the retained lines, namely the environments \texttt{lemma} on separate counters, together with the standard packages those lines need.  The retained environments print in this document as Lemma 1 in order of appearance, whereas the article prints the same items with the numbering of its own full document.  The complete native source of the article as distributed in the arXiv e-print for this exact version is bundled unchanged as \texttt{sources/197691/source0.tex}.
\begin{lemma}\label{eta_lemma}
	
	Suppose $\mathrm{char}(\BF_q) = p$, factor $q - 1 = \ell \cdot m$ where $\ell$ is an odd prime with $\ell \nmid m$. Set $o_1 = o_1(x, y) = \frac{q - 1}{\gcd\pt{\log_\varepsilon(x/y), q - 1}}$ where $x, y \in \BF_q^\times$ and $o_2 = o_2(z) = \frac{\log_\omega(z)}{\gcd\pt{\log_\omega(z), q + 1}}$ where $z \in \BK_q^\times$. Then we have
	\[\eta(g) = \left\{\begin{aligned}
		&(q + 1)\cdot\one\fk{{\mu_{m}(x)}},&&\text{ if } g = c_1(x);\\
		&\pt{{q}{p^{-1}} + 1}\cdot\one\fk{{\mu_{m}(x)}}, &&\text{ if } g = c_2(x);\\
		&\one\fk{\mu_m(x)} + \one\fk{\mu_m(y)} + \gcd{(\log_\varepsilon\pt{{x}/{y}}, q - 1)}\cdot\one\fk{\mu_{m}(x^{o_1})},&&\text{ if } g = c_3(x, y);\\
		&\gcd(\log_\omega(z), q + 1)\cdot\one\fk{\mu_{m}\pt{\varepsilon^{o_2}}},&&\text{ if } g = c_4(z).
	\end{aligned}\right.\]	
	
\end{lemma}

\begin{proof}
	
	The method for calculating $\eta$ is as follows: Note that $\dim_{\BC}V_{m, 0} = q + 1$ and for any $f \in V_{m, 0}$, $f$ is completely determined by $f(x_1), \ldots, f(x_{q + 1})$ where $\cb{x_1, \ldots, x_{q + 1}}$ is a $B$-transversal of $\GL_2(q)$. Fix $g \in \GL_2(q)$ and suppose that $g(\varphi) = \varphi$ where $\varphi \in V_{m, 0}$. Next we examine how many independent equations are needed to relate $f(x_1), \ldots, f(x_{q + 1})$, say $k$ equations, so $\eta(g) = q + 1 - k$. Because $\eta$ is a class function and $\GL_2(q)$ has four conjugacy classes, there are four cases to discuss. 
	
	\textbf{Case 1.} Suppose $\jz{x}{}{}{x}(\varphi) = \varphi$, then for any $t \in T_1$, we have
	\[\jk{t, \varphi}
	= \jk{t, \jz{x}{}{}{x}(\varphi)}
	= \jk{\jz{x}{}{}{x}t, \varphi}
	= \mu_{m}(x)\cdot\jk{t, \varphi}.\]
	Hence $\eta\pt{c_1(x)} = (q + 1)\cdot\one\fk{{\mu_{m}(x)}}$.
	
	\textbf{Case 2.} Suppose $\jz{x}{}{1}{x}(\varphi) = \varphi$, then for any $t \in \BF_q$, we have
	\[\begin{aligned}
		\jk{\jz{1}{}{t}{1}, \varphi}
		&= \jk{\jz{1}{}{t}{1}, \jz{x}{}{1}{x}(\varphi)}
		&&= \jk{\jz{1}{}{t}{1}\jz{x}{}{1}{x}, \varphi} \\
		&= \jk{\jz{x}{}{t\cdot x + 1}{x}, \varphi}
		&&= \jk{\jz{x}{}{}{x}\jz{1}{}{t + x^{-1}}{1}, \varphi} \\
		&= \mu_{m}(x)\cdot\jk{\jz{1}{}{t + x^{-1}}{1}, \varphi}
		&&= \mu_{m}(x^p)\cdot\jk{\jz{1}{}{t}{1}, \varphi}.
	\end{aligned}\]
	and
	\[\begin{aligned}
		\jk{\jz{}{1}{1}{}, \varphi} 
		&= \jk{\jz{}{1}{1}{}, \jz{x}{}{1}{x}(\varphi)}
		&= \jk{\jz{}{1}{1}{}\jz{x}{}{1}{x}, \varphi} \\
		&= \jk{\jz{1}{x}{x}{}, \varphi}
		&= \jk{\jz{x}{1}{}{x}\jz{}{1}{1}{}, \varphi} \\
		&= \mu_{m}(x) \cdot \jk{\jz{}{1}{1}{}, \varphi}.
	\end{aligned}\]
	Hence $\eta\pt{c_2(x)} = \frac{q}{p}\cdot\one\fk{{\mu_{m}(x^p)}} + \one\fk{{\mu_{m}(x)}} = \pt{q\cdot p^{-1} + 1}\cdot\one\fk{\mu_m(x)}$ since $p \nmid q - 1$.
	
	\textbf{Case 3.} Suppose $\jz{x}{}{}{y}(\varphi) = \varphi$ and the order of $x\cdot y^{-1}$ is $o_1$ in $\BF_q^\times$, explicitly, $o_1 = \frac{q - 1}{\gcd\pt{\log_\varepsilon(x/y), q - 1}}$. Then for any $t \in \BF_q$, we have
	\[\begin{aligned}
		\jk{\jz{1}{}{t}{1}, \varphi}
		&= \jk{\jz{1}{}{t}{1}, \jz{x}{}{}{y}(\varphi)} 
		&&= \jk{\jz{1}{}{t}{1}\jz{x}{}{}{y}, \varphi} \\
		&= \jk{\jz{x}{}{t\cdot x}{y}, \varphi} 
		&&= \jk{\jz{x}{}{}{y}\jz{1}{}{t\cdot x\cdot y^{-1}}{1}, \varphi} \\
		&= \mu_m(x)\cdot\jk{\jz{1}{}{t\cdot x\cdot y^{-1}}{1}, \varphi}  
		&&= \mu_m(x^{o_1})\cdot\jk{\jz{1}{}{t}{1}, \varphi}.
	\end{aligned}\]
	and
	\[\begin{aligned}
		\jk{\jz{}{1}{1}{}, \varphi}
		&= \jk{\jz{}{1}{1}{}, \jz{x}{}{}{y}(\varphi)}
		&= \jk{\jz{}{1}{1}{}\jz{x}{}{}{y}, \varphi} \\
		&= \jk{\jz{}{y}{x}{}, \varphi}
		&= \jk{\jz{y}{}{}{x}\jz{}{1}{1}{}, \varphi} \\
		&= \mu_m(y)\cdot\jk{\jz{}{1}{1}{}, \varphi}.
	\end{aligned}\]
	Hence $\eta\pt{c_3(x, y)} = \one\fk{\mu_m(x)} + \one\fk{\mu_m(y)} + \frac{q - 1}{o_1}\cdot\one\fk{\mu_{m}(x^{o_1})} = \one\fk{\mu_m(x)} + \one\fk{\mu_m(y)} + \gcd\pt{\log_\varepsilon(x/y), q - 1}\cdot\one\fk{\mu_m(x^{o_1})}$.
	
	\textbf{Case 4.} Suppose $\Theta(z)(\varphi) = \varphi$, then for any $0 \ls i \ls q$, we have (we omit the notation $\Theta$)
	\[\begin{aligned}
		\jk{\omega^i, \varphi}
		&= \jk{\omega^i, \omega^{\log_\omega(z)}(\varphi)}
		&&= \jk{\omega^i \cdot \omega^{\log_\omega(z)}, \varphi} \\
		&= \jk{\omega^{i + \log_\omega(z)}, \varphi} 
		&&= \jk{\omega^{i + \lcm(\log_\omega(z), q + 1)}, \varphi} \\
		&= \jk{\pt{\omega^{q + 1}}^{o_2}\cdot\omega^i, \varphi}
		&&= \jk{\varepsilon^{o_2}\cdot\omega^i, \varphi} \\
		&= \mu_{m}(\varepsilon^{o_2})\cdot\jk{\omega^i, \varphi}.
	\end{aligned}\]
	Hence $\eta(c_4(z)) = \gcd\pt{\log_\omega(z), q + 1}\cdot\one\fk{{\mu_{m}(\varepsilon^{o_2})}}$.\qedhere
	
\end{proof}

\end{document}
